\section{Embodiment and Experiential Evaluation}
\label{sec:introduction_embodiment_experience}

The trade-off between ergometer-congruent and boat-centric representation is especially relevant for virtual embodiment. Virtual embodiment concerns the extent to which users experience a virtual body or body part as related to their own body. Commonly discussed components include the sense of agency over the virtual body's actions, the sense of body ownership, and self-location in relation to the virtual body or space.
% TODO cite: sense of embodiment framework, e.g. agency, ownership, self-location.

The ergometer-congruent condition is closely related to principles from virtual hand and body ownership research. In such paradigms, embodiment can be influenced by whether seen and felt body positions, movements, or touches are spatially and temporally aligned. In the rowing case, the real and virtual hands or handle can potentially be placed in close correspondence: the user feels the physical handle and sees the virtual hands or handle move in the same way.
% TODO cite: rubber hand illusion / virtual hand illusion / visuomotor synchrony source.

However, rowing also differs from classic hand-ownership paradigms. It is an active, rhythmic, tool-mediated, whole-body exercise rather than a passive illusion involving a single limb. The user does not merely observe a virtual hand; they exert force, coordinate multiple body segments, and interact with a physical machine. This makes indoor rowing an interesting extension of embodiment questions into a sport and exercise context.

The boat-centric condition may be less directly congruent with the user's felt hand and handle movement, but it may better express what users understand as rowing. A virtual boat, oars, water, and forward movement can contribute to perceived sport realism and plausibility. For a rower or rowing-interested user, this representation may feel more meaningful or more appropriate to the activity, even if the mapping between real and virtual movement is less direct.

For this reason, the thesis does not treat embodiment as the only relevant outcome. The primary focus remains on embodiment-related experience, particularly agency, ownership, self-location, and overall sense of embodiment. However, the study also considers perceived rowing realism and user preference as design-facing measures. A representation can be sensorimotor-congruent but still be experienced as a poor depiction of rowing; conversely, a representation can be less congruent but preferred because it better captures the intended sport scenario.

Preference is therefore treated as an exploratory complement rather than as the main theoretical construct. It is included because practical XR sport systems are not evaluated only by whether they produce embodiment, but also by whether users find the representation convincing, appropriate, and desirable for the activity. At the same time, preference must be interpreted cautiously, because it may depend on visual polish, prior rowing experience, individual expectations, and details of the implementation.
